% Part: second-order-logic
% Chapter: syntax-and-semantics
% Section: satisfaction

\documentclass[../../../include/open-logic-section]{subfiles}

\begin{document}

\olfileid{sol}{syn}{sat}

\olsection{Relasi Nyukupi}

\begin{explain}
Kanggo netepake relasi nyukupi $\Sat{M}{!A}[s]$ kanggo
!!{formula}s orde kapindho, kita kudu njembarake definisi supaya
ngemot !!{variable}s orde kapindho. Pangerten !!a{structure}
ing logika orde kapindho padha kaya ing logika orde kapisan.
Bedane mung ing pangaji variabel~$s$: saiki pangaji iku kudu
menehi nilai ora mung marang !!{variable}s orde kapisan,
nanging uga marang !!{variable}s orde kapindho.
\end{explain}

\begin{defn}[Pangaji Variabel]
\emph{Pangaji variabel}~$s$ kanggo !!a{structure}~$\Struct{M}$
iku fungsi kang nggandhengake saben
\begin{enumerate}
\item !!{variable} objek~$\Obj{v_i}$ karo sawijining unsur
  saka~$\Domain M$, yaiku $s(\Obj{v_i}) \in \Domain{M}$
\item variabel relasi $n$-papan~$\Obj{V_i^n}$ karo relasi
  $n$-papan ing~$\Domain{M}$, yaiku
  $s(\Obj{V_i^n}) \subseteq \Domain{M}^n$;
\item variabel fungsi $n$-papan~$\Obj{u_i^n}$ karo fungsi
  $n$-papan saka $\Domain{M}$ menyang $\Domain{M}$, yaiku
  $s(\Obj{u_i^n})\colon \Domain{M}^n \to \Domain{M}$;
\end{enumerate}
\end{defn}

\begin{explain}
!!^a{structure} menehi !!a{value} marang saben !!{constant}
lan !!{function}, dene pangaji variabel orde kapindho menehi objek
lan fungsi marang saben variabel objek lan fungsi. Bareng-bareng,
loro-lorone ngidini kita menehi nilai marang saben term.
\end{explain}

\begin{defn}[\usetoken{S}{value} saka Term]
Yen $t$ iku term saka basa~$\Lang L$, $\Struct M$ iku
!!a{structure} kanggo~$\Lang L$, lan $s$ iku pangaji
!!a{variable} kanggo~$\Struct M$, mula \emph{!!{value}}~$\Value{t}{M}[s]$
ditetepake kaya kanggo term orde kapisan, ditambah klausa iki:
\begin{quote}
\indcase{t}{\Atom{u}{t_1, \ldots, t_n}}{
\[
\Value{\indfrm}{M}[s] = s(u)(\Value{t_1}{M}[s], \ldots,
\Value{t_n}{M}[s]).
\]}
\end{quote}
\end{defn}

\begin{defn}[Varian-$x$]
Yen $s$ iku pangaji !!a{variable} kanggo !!a{structure}~$\Struct M$,
mula saben pangaji !!{variable} $s'$ kanggo $\Struct M$ kang
bedane karo $s$ paling akeh mung ing nilai kang diwenehake marang
$x$ diarani \emph{varian-$x$} saka~$s$. Yen $s'$ iku varian-$x$
saka $s$, kita nulis $\varAssign{s'}{s}{x}$. (Semono uga kanggo
variabel orde kapindho $X$ utawa $u$.)
\end{defn}

\begin{defn}
  Yen $s$ iku pangaji !!a{variable} kanggo
  !!a{structure}~$\Struct M$ lan $m \in \Domain{M}$, mula
  pangaji~$\Subst{s}{m}{x}$ iku pangaji !!{variable} kang
  ditetepake dening
  \[\Subst{s}{m}{y} = \begin{cases}
    m & \text{yen } y \ident x\\
    s(y) & \text{yen ora},
  \end{cases}\]
  Yen $X$ iku !!{variable} relasi $n$-papan lan $M \subseteq
  \Domain{M}^n$, mula $\Subst{s}{M}{X}$ iku pangaji
  !!{variable} kang ditetepake dening
  \[\Subst{s}{M}{y} = \begin{cases}
    M & \text{yen } y \ident X\\
    s(y) & \text{yen ora}.
  \end{cases}\]
  Yen $u$ iku !!{variable} fungsi $n$-papan lan
  $f\colon \Domain{M}^n \to \Domain{M}$, mula
  $\Subst{s}{f}{u}$ iku pangaji !!{variable} kang ditetepake
  dening
  \[\Subst{s}{f}{y} = \begin{cases}
    f & \text{yen } y \ident u\\
    s(y) & \text{yen ora}.
  \end{cases}\]
  Ing saben kasus, $y$ bisa dadi !!{variable} orde kapisan utawa
  orde kapindho apa wae.
\end{defn}

\begin{defn}[Relasi Nyukupi]
Kanggo !!{formula}s orde kapindho~$!A$, definisi relasi nyukupi
kaya ing \olref[fol][syn][sat]{defn:satisfaction}, ditambah:
\begin{enumerate}
\item \indcase{!A}{\Atom{X^n}{t_1, \dots, t_n}}{$\Sat{M}{\indfrm}[s]$
  yen lan mung yen $\langle \Value{t_1}{M}[s], \dots, \Value{t_n}{M}[s] \rangle \in
  s(X^n)$.}
\tagitem{prvAll}{%
  \indcase{!A}{\lforall[X][!B]}{$\Sat{M}{\indfrm}[s]$ yen lan mung yen kanggo saben $M
  \subseteq \Domain{M}^n$, $\Sat{M}{!B}[\Subst{s}{M}{X}]$.}}{}

\tagitem{prvEx}{%
  \indcase{!A}{\lexists[X][!B]}{$\Sat{M}{\indfrm}[s]$ yen lan mung yen ana paling ora
  siji $M \subseteq \Domain{M}^n$ saengga
  $\Sat{M}{!B}[\Subst{s}{M}{X}]$.}}{}

\tagitem{prvAll}{%
  \indcase{!A}{\lforall[u][!B]}{$\Sat{M}{\indfrm}[s]$ yen lan mung yen kanggo saben
  $f\colon \Domain{M}^n \to \Domain{M}$,
  $\Sat{M}{!B}[\Subst{s}{f}{u}]$.}}{}

\tagitem{prvEx}{%
  \indcase{!A}{\lexists[u][!B]}{$\Sat{M}{\indfrm}[s]$ yen lan mung yen ana paling ora
  siji $f\colon \Domain{M}^n \to \Domain{M}$ saengga
  $\Sat{M}{!B}[\Subst{s}{f}{u}]$.}}{}
\end{enumerate}
\end{defn}

\begin{ex}
  Gatekna !!{formula}
  $\lforall[z][(\Atom{X}{z} \liff \lnot \Atom{Y}{z})]$.
  Formula iki ora ngemot kuantor orde kapindho, nanging ngemot
  !!{variable}s orde kapindho $X$ lan~$Y$ (ing kene dianggep
  siji-papan). !!{sentence} orde kapisan kang cocog,
  $\lforall[z][(\Atom{P}{z} \liff \lnot \Atom{R}{z})]$,
  nyatakake yen apa wae kang kalebu interpretasi~$P$ ora
  kalebu interpretasi~$R$, lan kosok balene. Ing
  !!a{structure}, interpretasi !!a{predicate}~$P$
  diwenehake dening interpretasi~$\Assign{P}{M}$. Nanging
  kanggo !!{variable}s orde kapindho kaya $X$ lan~$Y$,
  interpretasi diwenehake dudu dening !!{structure} kasebut
  dhewe, nanging dening pangaji !!a{variable}. Amarga
  !!{formula} orde kapindho iki dudu !!a{sentence} (formula iki
  ngemot !!{variable}s bebas $X$ lan~$Y$), relasi nyukupine
  mung bisa ditemtokake relatif marang
  !!a{structure}~$\Struct{M}$ bebarengan karo pangaji
  !!a{variable}~$s$.

  $\Sat{M}{\lforall[z][(\Atom{X}{z} \liff \lnot \Atom{Y}{z})]}[s]$
  bener yen saben !!{element}s ing domain kalebu ing~$s(X)$ persis yen
  !!{element}s iku ora kalebu ing~$s(Y)$, yaiku yen lan mung yen
  $s(Y) = \Domain{M} \setminus s(X)$. Cathetan koreksi sumber OLPL-803: mung ora irisan durung cukup; bikondisional iki mbutuhake Y persis komplemen X ing domain. Tuladhane, upamana
  $\Domain{M} = \{1, 2, 3\}$. Amarga ora ana !!{predicate}s,
  !!{function}s, utawa !!{constant}s kang melu, mung domain
  saka~$\Struct{M}$ kang wigati. Saiki, kanggo
  $s_1(X) = \{1, 2\}$ lan $s_1(Y) = \{3\}$, kita nduweni
  $\Sat{M}{\lforall[z][(\Atom{X}{z}
      \liff \lnot \Atom{Y}{z})]}[s_1]$.

  Kosok baline, yen $s_2(X) = \{1, 2\}$ lan
  $s_2(Y) = \{2, 3\}$, kita nduweni
  $\Sat/{M}{\lforall[z][(\Atom{X}{z} \liff \lnot
      \Atom{Y}{z})]}[s_2]$. Sebabe,
  $\Sat{M}{\Atom{X}{z}}[\Subst{s_2}{2}{z}]$ (amarga
  $2 \in \Subst{s_2}{2}{z}(X)$), nanging
  $\Sat/{M}{\lnot \Atom{Y}{z}}[\Subst{s_2}{2}{z}]$
  (amarga uga $2 \in \Subst{s_2}{2}{z}(Y)$).
\end{ex}

\begin{ex}
  $\Sat{M}{\lexists[Y][(\lexists[y][\Atom{Y}{y}] \land
  \lforall[z][(\Atom{X}{z} \liff \lnot \Atom{Y}{z})])]}[s]$
  bener yen ana $N \subseteq \Domain{M}$ saengga
  $\Sat{M}{(\lexists[y][\Atom{Y}{y}] \land \lforall[z][(\Atom{X}{z}
  \liff \lnot \Atom{Y}{z})])}[\Subst{s}{N}{Y}]$. Iki kedadeyan
  yen $N \neq \emptyset$ (saengga
  $\Sat{M}{\lexists[y][\Atom{Y}{y}]}[\Subst{s}{N}{Y}]$) lan,
  kaya ing tuladha sadurunge, $N = \Domain{M} \setminus s(X)$.
  Kanthi tembung liya,
  $\Sat{M}{\lexists[Y][(\lexists[y][\Atom{Y}{y}] \land
  \lforall[z][(\Atom{X}{z} \liff \lnot \Atom{Y}{z})])]}[s]$
  yen lan mung yen $\Domain{M} \setminus s(X)$ ora kosong,
  yaiku $s(X) \neq \Domain{M}$. Dadi, !!{formula} iki
  dicukupi, upamane, yen $\Domain{M} = \{1, 2, 3\}$
  lan $s(X) = \{1, 2\}$, nanging ora yen
  $s(X) = \{1, 2, 3\} = \Domain{M}$.

  Amarga !!{formula} iki ora dicukupi yen
  $s(X) = \Domain{M}$, mula !!{sentence}
  \[
  \lforall[X][\lexists[Y][(\lexists[y][\Atom{Y}{y}] \land
      \lforall[z][(\Atom{X}{z} \liff \lnot \Atom{Y}{z})])]]
  \]
  ora tau dicukupi: kanggo saben !!{structure}~$\Struct{M}$,
  pangaji~$s(X) = \Domain{M}$ bakal ndadekake !!{sentence}
  kasebut salah. Kosok baline, ukara
  \[
  \lexists[X][\lexists[Y][(\lexists[y][\Atom{Y}{y}] \land
      \lforall[z][(\Atom{X}{z} \liff \lnot \Atom{Y}{z})])]]
  \]
  dicukupi relatif marang pangaji~$s$ apa wae, amarga kita
  tansah bisa nemokake $M \subseteq \Domain{M}$ kanthi
  $M \neq \Domain{M}$ (umpamane, $M = \emptyset$).
\end{ex}

\begin{ex}
  !!{sentence} orde kapindho~$\lforall[X][\lforall[y][X(y)]]$
  nyatakake yen saben relasi $1$-papan, yaiku saben sipat,
  nyakup saben objek. Pratelan iki cetha ora tau bener,
  amarga ing saben~$\Struct{M}$, kanggo pangaji variabel~$s$
  kanthi $s(X) = \emptyset$ lan $s(y) = a \in \Domain{M}$,
  kita nduweni $\Sat/{M}{X(y)}[s]$. Tegese,
  $!A \lif \lforall[X][\lforall[y][X(y)]]$ iku ekuivalen
  ing logika orde kapindho karo $\lnot !A$, yaiku:
  $\Sat{M}{!A \lif \lforall[X][\lforall[y][X(y)]]}$
  yen lan mung yen $\Sat{M}{\lnot !A}$. Kanthi tembung
  liya, ing logika orde kapindho kita bisa netepake $\lnot$
  nganggo $\lforall$ lan~$\lif$.
\end{ex}

\begin{prob}
  Tuduhna yen ing logika orde kapindho $\lforall$ lan~$\lif$
  bisa netepake konektif liyane:
  \begin{enumerate}
    \item Buktèkna yen ing logika orde kapindho $!A \land !B$
    ekuivalen karo $\lforall[X][(!A \lif (!B \lif \lforall[x][X(x)])\lif
    \lforall[x][X(x)])]$.
    \item Temokna formula orde kapindho kang mung nggunakake
    $\lforall$ lan $\lif$ lan ekuivalen karo $!A \lor !B$.
  \end{enumerate}
\end{prob}

\end{document}
